\section{模型假设与符号说明}
\subsection{模型假设及其用途}
\begin{enumerate}[label=（\arabic*）]
\item 药材横截面内满足轴对称，材料在周向上均匀，表面交换条件沿周向一致。忽略轴向梯度和端面交换，建立长度中部截面的一维径向模型。该假设用于将传递方程降为径向形式；长时间计算的端面影响作为局限保留。
\item 初始温度与干基含水率均匀；问题一至三半径与长度固定。问题四长度仍取0.25 m，骨架按给定$R(t)$作比例径向收缩，干物质不流失。比例收缩用于由外半径轨迹闭合内部速度场。
\item 题给$\rho c_p$作为有效体积热容量，采用Fourier导热和有效Fick扩散描述内部传递。热方程不显式包含汽化潜热、水分携焓及机械功。此处理用于在题给数据下封闭模型，不代表这些机制已被实验排除。
\item 沿用$h=25\ \mathrm{W/(m^2\cdot K)}$和$h_m=8\times10^{-7}$ m/s作为表面交换系数。将$h_m(C_s-C_a)$理解为题给信息下的等效水分交换关系。气相与固相的kg/kg具有不同干基载体，未将$C_a$直接认定为实测药材平衡含水率。
\item 固定域采用均匀、恒定的参考干物质密度；收缩域按干物质守恒更新该密度。它与用于热方程的经验$\rho(C)$分别记账，不强制二者逐点满足湿密度换算关系。
\end{enumerate}

径向近似在问题一中可作尺度说明。初始长度与直径之比为6.25，两端总面积与侧面积之比为0.08。附录2给出热扩散率$\alpha=0.36/(820\times2600)=1.6886\times10^{-7}$ m$^2$/s，1800 s内的传播尺度$\sqrt{\alpha t}=0.0174$ m，小于中截面到端面的距离0.125 m。这支持短时中截面的径向近似，但面积比和传播尺度都不是后续数日预测的误差上界。

\subsection{主要符号}
\begin{center}\small
\renewcommand{\arraystretch}{1.13}
\begin{tabular}{p{2.4cm}p{7.5cm}>{\raggedright\arraybackslash}p{3cm}}
\toprule 符号 & 含义 & 单位 \\\midrule
$t,r$ & 入炉后的时间、当前径向距离 & s，m \\
$R,R(t),L$ & 固定半径、随时间变化的半径、长度 & m \\
$T,T_a$ & 药材局部温度、烘房温度 & $^\circ$C \\
$T_K$ & 药材局部绝对温度，$T_K=T+273.15$ & K \\
$C,C_a$ & 药材干基含水率、烘房题给水分量 & kg/kg \\
$\rho,c_p,k$ & 有效密度、比热容、导热系数 & kg/m$^3$，J/(kg$\cdot$K)，W/(m$\cdot$K) \\
$B=\rho c_p,D$ & 有效体积热容量、水分扩散系数 & J/(m$^3\cdot$K)，m$^2$/s \\
$h,h_m$ & 对流换热系数、有效传质系数 & W/(m$^2\cdot$K)，m/s \\
$s,v$ & 当前体积干物质密度、骨架径向速度 & kg/m$^3$，m/s \\
$\xi,U,\Theta$ & 材料坐标、该坐标下的含水率和温度 & 1，kg/kg，$^\circ$C \\
$W_i,w_i$ & 物理域、材料域控制体的径向权重 & m$^2$，1 \\
$M,t_*,t_{\mathrm{rep}}$ & 全域最大含水率、临界时刻、报告时刻 & kg/kg，s，s \\
\bottomrule
\end{tabular}\end{center}
下标$s$表示当前表面，$i$表示内部计算节点。正文中的时间公式以s为单位，时长结果及表3至表6换算为h。
